\begin{frame}[label=q5e2]{5e.2 · Vistas temporales: dos JOIN de perfil}
  \small
  \pregunta{¿Qué combinación de grupo etario y sexo concentra más personas y qué porcentaje del total representa?}
  \entrada{Vistas \texttt{denuncias}, \texttt{dim\_edad}, \texttt{dim\_sexo}. Claves: \texttt{rango\_edad} y \texttt{sexo}.}
  \flujo{Vistas registradas\\En esquema.scala}{\texttt{JOIN dim\_edad}\\\texttt{JOIN dim\_sexo}}{\texttt{GROUP BY} grupo y sexo\\Suma y \% nacional}{18 combinaciones\\Orden por personas}
  \begin{columns}[T,onlytextwidth]
    \column{0.55\textwidth}
    \textbf{Código Scala original (fragmento)}
    \codigo{sql/05e_2_grupoetario_sexo}{17}{17}
    \codigo{sql/05e_2_grupoetario_sexo}{22}{26}
    \column{0.42\textwidth}
    \textbf{Salida}\\[5pt]
    {\footnotesize
\begin{tabular}{@{}llll@{}}
\toprule
Grupo & Sexo & Personas & \% país\\
\midrule
Adolescente & F & 51\,021 & 42.02\\
Adolescente & M & 13\,309 & 10.96\\
Adulto & M & 12\,389 & 10.20\\
\bottomrule
\end{tabular}
    }
  \end{columns}
  \respuesta{Adolescentes mujeres reúnen 51\,021 personas: $100\times51\,021/121\,427=42.02\%$. Es una participación en las denuncias, no una probabilidad individual de desaparición.}
  \origen{sql}{05e_2_grupoetario_sexo}
  \note{30 s. Seguir ambos joins hasta el porcentaje. Las vistas se registraron con createOrReplaceTempView en esquema.scala.}
\end{frame}
